Für die härtere Feder gilt $D_\text{h} = (1+p)\cdot D_\text{w}$ mit $p = \pHO$. Bei gleicher Kraft $F$ ist
\begin{align}
 x_\text{w} &= \frac{F}{D_\text{w}}, \qquad x_\text{h} = \frac{F}{D_\text{h}} = \frac{F}{(1+p)\,D_\text{w}}
 \quad\Rightarrow\quad \frac{x_\text{h}}{x_\text{w}} = \frac{1}{1+p}
\end{align}
Die relative Änderung der Ausdehnung ist damit
\begin{align}
 \frac{x_\text{h}-x_\text{w}}{x_\text{w}} &= \pxF \\
   &= \frac{1}{1+\pHO}-1 \\
   &= \px \approx \result{\pxP{0}{2}}
\end{align}
Die härtere Feder dehnt sich also um gut einen Viertel weniger aus (und nicht um \pHO, weil die Ausdehnung umgekehrt proportional zur Federkonstante ist).
